\documentclass[11pt]{amsart}

\usepackage{amsmath,amssymb}
\usepackage{hyperref}
\usepackage{orcidlink}

\newtheorem{theorem}{Theorem}[section]
\newtheorem{lemma}[theorem]{Lemma}
\newtheorem{proposition}[theorem]{Proposition}
\newtheorem{corollary}[theorem]{Corollary}
\theoremstyle{definition}
\newtheorem{definition}[theorem]{Definition}
\newtheorem{example}[theorem]{Example}
\theoremstyle{remark}
\newtheorem{remark}[theorem]{Remark}

\title{A Formalization of the Elligator Family of Uniform Encodings
  to Elliptic Curves}

\author{Chris Anto Fröschl\orcidlink{0009-0006-3562-5395}}
\address{Munich University of Applied Sciences, Department of
  Computer Science and Mathematics, Germany}
\email{chris.froeschl@hm.edu}

\author{Prof. Matthias Güdemann\orcidlink{0000-0002-1002-6023}}
\address{Munich University of Applied Sciences, Department of
  Computer Science and Mathematics, Germany}
\email{matthias.guedemann@hm.edu}

\subjclass[2020]{03B70, 68V20, 94A60, 11T71, 14G50}
\keywords{Elligator, elliptic curves, elliptic-curve cryptography,
  Curve1174, uniform random encoding, formal verification, Lean}

\begin{document}

\begin{abstract}
First Elligator formalization (what/who), What formalized (Ell1/2, Legendre), where stored (CSLib, what?), advances (prime power, computability, test vector testbenching...)
\end{abstract}

\maketitle

\section{Introduction}
\label{sec:introduction}

Elliptic-curve cryptography routinely transmits points on a fixed
public curve, and such points are easily distinguished from uniform
random bit strings: an observer who can tell elliptic-curve traffic
apart from random traffic can selectively block it. Bernstein,
Hamburg, Krasnova and Lange's Elligator~\cite{bernstein2013a}
addresses this by giving efficient, invertible encodings between
(most of) the points of a suitable elliptic curve and uniform bit
strings, so that a random curve point can be transmitted as what
looks like a random string. The paper introduces two such encodings,
Elligator~1 and Elligator~2, together with Curve1174, a concrete
Edwards curve suited to Elligator~1. Later work has extended the
approach to further curve shapes, including a Legendre-curve variant
due to Saha and Karati~\cite{gourab2026a}.

Despite substantial practical uptake --- Elligator has been
numerically tested~\cite{elligatorcryptotests} and implemented
across a range of libraries and protocols~\cite{elligatororg} --- the
underlying theory has, to our knowledge, never been formalized in a
proof assistant. This paper reports on such a formalization, carried
out in CSLib, a Lean~4 library that currently has essentially no
cryptographic primitives, not even basic elliptic-curve definitions.
Elligator is accordingly the first substantial elliptic-curve
encoding family to enter CSLib.

\paragraph{Contributions.}
\begin{enumerate}
  \item A formalization of elliptic curves and the curve primitives
    needed to state Elligator, none of which previously existed in
    CSLib.
  \item A formalization of Elligator~1 and Curve1174, and a
    generalization of the viability lemmas underlying Elligator~1
    from prime fields to prime power fields (Section~\ref{sec:elligator1}).
  \item A formalization of Elligator~2 and of the Legendre-curve
    variant of Saha and Karati as instances of a common encoding
    family alongside Elligator~1 (Section~\ref{sec:elligator2-legendre}).
  \item Computable, extractable definitions throughout, so that the
    formalized encodings are not merely existence statements but can
    be run and checked against published test
    vectors (Section~\ref{sec:computability}).
\end{enumerate}

\paragraph{Outline.}
Section~\ref{sec:background} recalls the curve and field background.
Section~\ref{sec:elligator1} covers Elligator~1 and its
generalization. Section~\ref{sec:elligator2-legendre} covers
Elligator~2 and the Legendre-curve variant.
Section~\ref{sec:curve1174} covers Curve1174.
Section~\ref{sec:computability} discusses computability and
validation against test vectors. Section~\ref{sec:related-work}
discusses related work and Section~\ref{sec:conclusion} concludes.


\section{Background}
\label{sec:background}

This section fixes the curve and field notions used throughout, all
newly formalized in CSLib for this work unless noted otherwise.

state the actual CSLib definitions of finite fields,
quadratic residues/the Legendre symbol (or point at Mathlib's, if
reused), and the curve shapes used below (Edwards, Montgomery,
Legendre), matching the Lean source exactly.

\subsection{Finite fields and quadratic residues}
\label{subsec:fields}

\subsection{Curve shapes}
\label{subsec:curve-shapes}

definitions for the (twisted) Edwards curves used by
Elligator 1, the Montgomery curves used by Elligator 2, and the
Legendre curves $E_\lambda : y^2 = x(x-1)(x-\lambda)$ used by the
Legendre-curve variant.


\section{Elligator 1 and its generalization to prime powers}
\label{sec:elligator1}

state Bernstein et al.'s original viability conditions
and lemmas, then generalized statements side by side, and give
a proof sketch of the generalization.

Elligator~1~\cite[\S3]{bernstein2013a} encodes bit strings as
points on a complete Edwards curve $x^2 + y^2 = 1 + dx^2y^2$ over a
prime field $\mathbb{F}_p$ with $p \equiv 3 \pmod 4$, for $d$ of a
specific non-square form. The viability of the construction --- that
the encoding is well-defined and injective on the intended domain ---
is established in the original paper for prime $p$.

\subsection{Generalizing to prime powers}
\label{subsec:prime-power-generalization}

We show that the viability lemmas underlying Elligator~1 do not in
fact require $p$ to be prime, only that the underlying field has
$q = p^k$ elements with $q \equiv 3 \pmod 4$: that is, the
construction is viable over any finite field of order a prime power
congruent to $3 \bmod 4$, not only prime fields. This is, to our
knowledge, a genuine extension of Bernstein et al.'s result rather
than a restatement of it.

\begin{theorem}[Generalized Elligator 1 viability]
\label{thm:elligator1-prime-power}
TODO placeholder
Let $q = p^k$ be a prime power with $q \equiv 3 \pmod 4$. Then the
Elligator~1 encoding is well-defined and injective on
[TODO: the intended domain] over $\mathbb{F}_q$.
\end{theorem}


\section{Elligator 2 and the Legendre-curve variant}
\label{sec:elligator2-legendre}

We formalize both of these alongside Elligator~1, collecting the
three under a common presentation in CSLib.

\subsection{Elligator 2}
\label{subsec:elligator2}

TODO formal statement and theorem.

\subsection{The Legendre-curve variant}
\label{subsec:legendre-variant}

TODO formal statement and theorem.

\subsection{A common family}
\label{subsec:common-family}

TODO if there is a genuine shared interface/theorem here
(e.g. a common correctness criterion all three encodings satisfy),
state it. If not, at minimum give a comparison table of domain,
codomain, applicable curve shapes, and injectivity radius across
the three encodings.


\section{Curve1174}
\label{sec:curve1174}

TODO state the concrete curve equation/parameters, and
which properties of it (completeness, group order, twist security,
suitability for Elligator 1) are formalized vs. assumed/imported.

Curve1174~\cite[\S4]{bernstein2013a} is the concrete Edwards curve
$x^2 + y^2 = 1 - 1174x^2y^2$ over $\mathbb{F}_p$ with
$p = 2^{251} - 9$, introduced alongside Elligator~1 as a curve
satisfying the conditions required for the encoding together with
the security properties expected of a cryptographic curve. We
formalize its definition and the properties needed to instantiate
the Elligator~1 encoding of Section~\ref{sec:elligator1} on it.


\section{Computability and validation}
\label{sec:computability}

TODO describe the extraction/execution mechanism used
(Lean's eval, code generation, or otherwise), and report the
actual validation.

A definition that merely proves the existence of an encoding is of
limited practical use; we give computable, extractable definitions
for every encoding in this paper, so that they produce actual byte
sequences that can be compared against independent implementations.
This matters for CSLib's stated goal of being practically reusable,
not only mathematically faithful.

\subsection{Extraction}
\label{subsec:extraction}

TODO how the Lean definitions are made executable, and any
performance caveats.

\subsection{Validation against published test vectors}
\label{subsec:validation}

TODO report which test vectors were checked (e.g. from elligator.cr.yp.to or elligator.org) and the outcome.

We cross-checked the extracted implementations of
Section~\ref{subsec:extraction} against the test vectors published
alongside the original Elligator paper~\cite{elligatorcryptotests}
and the reference implementations catalogued at~\cite{elligatororg}.


\section{Related work}
\label{sec:related-work}

The Elligator constructions have been extended and adapted in
several directions we do not formalize here, including Elligator
Squared~\cite{tibouchi2014a}, which supports curves of
prime order not covered by the original constructions, and
Elligator-T~\cite{gourab2026a} for twisted Edwards curves. To
our knowledge, none of the Elligator family has previously been
formalized in a proof assistant: existing validation has been
numerical~\cite{elligatorcryptotests} or by way of independent
implementations~\cite{elligatororg}, rather than machine-checked
proof.

CSLib itself has, prior to this work, essentially no cryptographic
primitives; a separate, unrelated CSLib effort formalizes binding
representations for the untyped lambda calculus, sharing this
paper's goal of building reusable, well-reviewed foundations rather
than one-off developments, but is otherwise unconnected in subject
matter.

\section{Conclusion}
\label{sec:conclusion}

We have given the first formalization of the Elligator family of
uniform encodings to elliptic curves, covering Elligator~1 (with a
generalization of its viability lemmas from prime to prime power
fields), Elligator~2, the Legendre-curve variant, and Curve1174,
together with computable definitions validated against published
test vectors. This is CSLib's first substantial cryptographic
development, and we hope it serves as a starting point for further
elliptic-curve formalization within the library.



\bibliographystyle{alpha}
\bibliography{refs}

\end{document}
